\section{Related Work}
\paragraph{Attacks.} FGSM~\cite{goodfellow2015fgsm} takes one signed-gradient step; PGD~\cite{madry2018pgd} iterates it inside an $\ell_\infty$ ball and is the standard strong first-order attack. DeepFool~\cite{moosavi2016deepfool} linearizes the decision boundary to find a minimal perturbation, and Carlini--Wagner (CW)~\cite{carlini2017cw} optimizes a margin loss with a norm penalty. Square Attack~\cite{andriushchenko2020square} is a score-based black-box attack using random search. BPDA~\cite{athalye2018obfuscated} replaces a non-differentiable defense by a differentiable approximation in the backward pass and showed that many preprocessing defenses give a false sense of security.

\paragraph{Preprocessing defenses.} JPEG compression~\cite{dziugaite2016jpeg}, its adaptations~\cite{das2018shield,liu2019featdist}, smoothing and bit-depth reduction all remove high-frequency perturbation energy. Feature squeezing~\cite{xu2018squeeze} turns this into a detector by comparing a model's output on the original and squeezed inputs. Feature distillation~\cite{liu2019featdist} reports that quality factors must be very low to resist BPDA, which hurts clean accuracy.

\paragraph{Face recognition.} Verification models such as FaceNet~\cite{schroff2015facenet} learn embeddings compared by distance. Adversarial vulnerability and simple perturbation robustness of face recognition were studied in~\cite{goswami2018unravelling}, and detection of attacks using deep feature and similarity measurements in~\cite{massoli2021detection}.

\paragraph{Gap (to the best of our limited search).} We did not find prior work that evaluates a pre/post-transformation \emph{embedding-consistency} detector for verification pairs with a Bayes decision layer against adaptive attackers. Our literature search was short and web-based, so this is a statement about what we found, not a proof of novelty.
